% ======================================================================
% Section 5 — The witness archive
% ======================================================================

\section{The witness archive}\label{sec:witness}

The archive has four shelves: exact refutations of the covering-radius
form (each a single rational point, checkable by one evaluation of
\eqref{eq:runnerform}, the sharpest carrying a complete two-sided
certificate---the exact covering radius itself), the arithmetic law
that governs the denominators of the
exact values (\S\ref{sec:weights}), the deepest-hole certification
ladder for the family $(1,\dots,n{-}1,m)$, and the branch-and-bound
certificate method together with its negative-control validation. Every
entry is re-asserted by two independent exact implementations in the
provenance run of Appendix~\ref{sec:prov}.

\subsection{Exact refutations of the covering-radius form}

Table~\ref{tab:refute} lists the witnesses. Each witness $c$ is a
rational point of the quotient torus in $c$-coordinates; the claimed
value is $D(c)$ from \eqref{eq:runnerform}, computed exactly and
independently by the linear-form/lattice-search implementation of the
program's records. The first entry carries the strongest certification
in the archive: the value $16/47$ is the covering radius itself,
certified from both sides---the witness below, and the branch-and-bound
upper certificate of \S\ref{sec:bb}, whose tree empties at $r=16/47$.
The second entry is \emph{argmax-certified}: the displayed value is
not merely attained at the witness point but is the maximum of $D$ on
the whole argmax neighborhood, certified exactly by the arrangement
analysis of \S\ref{sec:bb} (all vertices, ridges, and corners of the
region evaluated in rational arithmetic), and the argmax carries the
binding structure of \S\ref{sec:weights} that forces the denominator;
its global equality remains open (\S\ref{sec:bb}).
The harmonic instance is extremal for the $\delta$-form
($\delta=\mathrm{thr}_n$ exactly), so the value $16/47$ kills the
$\rho$-form at the very instance where the surviving equivalent form is
tight; the rung-2 instance $v=(1,2,3,4,5,12)$ was previously a
boundary row (witness value exactly $\mathrm{thr}_6$) and is now a
strict refutation; the $m=9$ row is upgraded from boundary to strict
witness; the $v=(1,2,3,4,5,6)$ and $v=(1,2,3,4,5,7)$ rows are strict
witnesses of long standing. The witnesses are non-torsion in every
case: the deepest holes are not the center classes.

\begin{table}[htbp]
\centering
\footnotesize
\caption{Exact witnesses against the covering-radius form
$\rho(V)\le\mathrm{thr}_n$. All values are exact rational lower bounds
for $\rho$; the first is \emph{exact-certified} (the branch-and-bound
tree empties at $r$ equal to the displayed value, so the value is the
covering radius itself), the second \emph{argmax-certified} (the
maximum of $D$ on the whole argmax neighborhood equals the displayed
value, exact arrangement analysis); both are depths of the deepest
hole, not merely point depths. Witness points are in $c$-coordinates
(\S\ref{sec:prelim}). ``$>\mathrm{thr}$?'' marks whether the value
\emph{strictly} exceeds the threshold: only strict rows falsify the
$\rho$-form; boundary rows (equality) separate the deepest hole from
the center class without falsifying it. The $\delta$ entries for
non-harmonic instances are computed values, not harmonic identities.}
\label{tab:refute}
\begin{tabularx}{\textwidth}{@{} l l l >{\raggedright\arraybackslash}X l c @{}}
\toprule
$V$ & $\delta=\tfrac12-\lambda$ & $\mathrm{thr}_n$ & witness $c$ &
$\rho$ & $>\mathrm{thr}$? \\
\midrule
$(1,2,3,4,5)$ & $1/3$ & $1/3$ &
$(\tfrac{122}{1175},\tfrac{688}{1175},\tfrac{284}{1175},
\tfrac{80}{1175})$ &
$16/47=.3404$ & yes (exact) \\
$(1,2,3,4,5,12)$ & $9/26$ & $5/14$ &
$(\tfrac{4714}{5875},\tfrac18,\tfrac58,\tfrac6{25},
\tfrac{2034}{5875})$ &
$17/47=.3617$ & yes (argmax) \\
$(1,2,3,4,5,9)$ & $5/14$ & $5/14$ &
$(\tfrac1{32},\tfrac{15}{32},\tfrac{1383}{10000},
\tfrac{32921}{40000},\tfrac{95473}{200000})$ &
$1793/5000=.3586$ & yes \\
$(1,2,3,4,5,6)$ & $5/14$ & $5/14$ &
$(0,\tfrac5{16},\tfrac{15}{16},\tfrac34,\tfrac12)$ & $29/80=.3625$ &
yes \\
$(1,2,3,4,5,7)$ & $1/3$ & $5/14$ &
$(\tfrac1{16},\tfrac9{16},0,\tfrac38,0)$ & $4/11{\approx}.3636$ & yes
\\
\bottomrule
\end{tabularx}
\end{table}

The $m=9$ witness deserves its sentence: the instance sits at the
threshold from the center side ($\delta=5/14=\mathrm{thr}_6$ computed),
so the boundary witness at exactly $5/14$ could not decide the
$\rho$-form either way; the strict witness $1793/5000>5/14$ does. What
it does not decide is the value itself: the exact covering radius at
$v=(1,2,3,4,5,9)$ is unknown. The recorded lower bound is marginally
deeper than the clean witness---an exactly evaluated dyadic search
point at depth $0.358639\ldots$ persists in the provenance data
(Appendix~\ref{sec:prov})---while no nontrivial upper certificate is
on record; the row's status is strict refutation, exact value unknown.

Three further exact families fill out the record---but it matters what
each refutes. The $n=5$ family $(1,2,3,4,m)$ refutes the
\emph{modular deepest-hole law}, not the $\rho$-form: the law's equality
clause (deepest hole at the center when $5\mid m$) fails exactly at the
multiples $m=5$ (now with the exact witness $\rho\ge16/47\ne\delta=1/3$)
and $m=15$ ($6/19>5/16=\delta$), while the deepest-hole-at-center
property fails at all twenty non-multiples (exact witnesses, the largest
margins $1/3>3/10$ at $m=6,7,8$). The $n=6$ family extinguishes the
modular law at every tested $m\in\{6,\dots,11,18,24\}$ (all multiples of
$6$ through $4n$ included; the narrowest modular margin is $7/5200$ at
$m=24$, witness value $71/208$ against $\delta=17/50$); against the
$\rho$-form itself those witnesses fall short of the threshold
($6/19<1/3$; $71/208<5/14$) and are recorded as deepest-hole evidence,
not covering-radius refutations. At $n=6$, $m=9$ the strict witness of
Table~\ref{tab:refute} upgrades the earlier boundary row; at $m=12$ the
witness $17/47$ upgrades another. A useful way to read the table:
the $\rho$-form is falsified at the extremal harmonic instance by the
exact argmax value $16/47$, at a second $n=6$ instance by the exact
argmax value $17/47$, and at
two further $n=6$ instances by exact strict witnesses; the
classification side (which instances have their deepest hole at the
center) is a different question, answered by the ladder below.

\subsection{The arithmetic of the deepest holes: a weight identity}
\label{sec:weights}

The exact values $16/47$ and $17/47$ are not accidents of the search:
their denominators are forced by a one-line linear identity among the
norm's defining forms, and the same identity explains why the exact
covering radii of these instances are rationals with small structured
denominators rather than generic reals. This subsection records the
identity, the two binding archetypes that activate it, and the resulting
denominator laws; everything is elementary and fully proved.

Recall from \S\ref{sec:prelim} that the quotient norm is the maximum of
the pair forms $\phi^{uv}$, and write the \emph{pair character}
$\chi_{uv}=(u{+}v)\,\phi^{uv}$, i.e.\ $\chi_{uv}(x)=v\,x_u-u\,x_v$ on
any lift $x$: a circle-valued linear map on the quotient torus
$\mathcal{X}_v/\Lambda$ (integer shifts of $x$ change $\chi_{uv}$ by
integers). For a lattice translate $z$, write
$k_{uv}(z)=v\,z_u-u\,z_v\in\Z$, so that
$\phi^{uv}(y-z)=\bigl(\chi_{uv}(y)-k_{uv}(z)\bigr)/(u{+}v)$.

\begin{lemma}[pair-character identity]\label{lem:chi}
For any three distinct speeds $a,b,c\in V$,
\begin{equation}\label{eq:chi}
c\,\chi_{ab}\;+\;a\,\chi_{bc}\;=\;b\,\chi_{ac}
\end{equation}
identically on the quotient torus.
\end{lemma}

\begin{proof}
Substitute $\chi_{uv}=v\,x_u-u\,x_v$ (any lift): the left side is
$c(bx_a-ax_b)+a(cx_b-bx_c)=bc\,x_a-ac\,x_b+ac\,x_b-ab\,x_c
=b(cx_a-ax_c)$, which is the right side. Every term is a function of
the quotient class, so the identity descends.
\end{proof}

Multiplying \eqref{eq:chi} through by the pair sums recovers the
\emph{weighted form} in terms of the pair forms themselves:
\begin{equation}\label{eq:weights}
c\,(a{+}b)\,\phi^{ab}\;+\;a\,(b{+}c)\,\phi^{bc}\;=\;
b\,(a{+}c)\,\phi^{ac}.
\end{equation}
The weight of each pair is the third speed times the pair sum. Three
pair forms of one speed triple are therefore never independent, and at a
point where all three are simultaneously \emph{binding}---each equal to
$\pm\rho$ in absolute value relative to its own nearest translate---the
identity becomes an arithmetic constraint on $\rho$ itself.

\begin{corollary}[denominator laws]\label{cor:denom}
Let $y$ be a point of the quotient torus and $\rho>0$.
\emph{(i) Triple binding.} If the three pair forms of a speed triple
$a<b<c$ are simultaneously binding at $y$,
$\phi^{uv}(y-z_{uv})=s_{uv}\,\rho$ with signs $s_{uv}\in\{\pm1\}$ and
translates $z_{uv}$, then
\[
\bigl[\,c\,s_{ab}(a{+}b)\;+\;a\,s_{bc}(b{+}c)\;-\;b\,s_{ac}(a{+}c)\,
\bigr]\,\rho\;\in\;\Z,
\]
the integer being $b\,k_{ac}-c\,k_{ab}-a\,k_{bc}$. In the sign pattern
$(+,+,-)$ and its negation---the pattern that occurs at both closed
argmaxes below---the bracket equals $2(ab{+}ac{+}bc)$, the sum of the
three weights in \eqref{eq:weights}, and
\[
\rho\;=\;\frac{\bigl|\,b\,k_{ac}-c\,k_{ab}-a\,k_{bc}\,\bigr|}
{2\,(ab+ac+bc)}.
\]
In particular the reduced denominator of $\rho$ divides
$2(ab+ac+bc)$.
\emph{(ii) Pair bisection.} If a single pair form $(u,v)$ binds at two
translates with opposite signs, $\phi^{uv}(y-z_1)=+\rho$ and
$\phi^{uv}(y-z_2)=-\rho$, then $2(u{+}v)\,\rho=k_{uv}(z_2)-k_{uv}(z_1)
\in\Z$: the reduced denominator of $\rho$ divides $2(u+v)$.
\end{corollary}

\begin{proof}
Part (i): binding means
$\chi_{uv}(y)=k_{uv}(z_{uv})+s_{uv}(u{+}v)\rho$ for each of the three
pairs; substitute the three expressions into \eqref{eq:chi} and collect
the $\rho$-terms. In the stated sign pattern the coefficient of $\rho$
is $c(a+b)+a(b+c)+b(a+c)=2(ab+ac+bc)$. Part (ii) is the same
substitution for the single identity $\chi_{uv}(y)-\chi_{uv}(y)=0$ with
the two translates: $k_{uv}(z_1)+(u{+}v)\rho=k_{uv}(z_2)-(u{+}v)\rho$.
\end{proof}

The scope of Corollary~\ref{cor:denom}(i) is worth stating
precisely, because the harmonic example exceeds it. The hypothesis is
that the three pair forms of \emph{one} speed triple bind---nothing is
assumed about the other pair forms, which may bind or not as they
please. An argmax with four or more binding forms therefore falls
under the corollary once for each speed triple whose three pair forms
all bind: the identity is applied to a three-element subset of the
binding set, and additional binding forms neither disturb nor dilute
it---every such triple with the $(+,+,-)$ sign pattern forces
$2(ab{+}ac{+}bc)\,\rho\in\Z$ independently. The harmonic argmax below
binds four pair forms; the triple $\{3,4,5\}$ among them forces
$2\cdot47\,\rho\in\Z$, and the fourth binding form (the pair $(2,5)$)
is extra structure the computation does not need. No comparable
constraint exists at a point where fewer than three pair forms bind in
the pattern; the $m=9$ witness below is a chain of dips that never
fully ties, and its clean value reflects the parametrized scan that
found it rather than any forced arithmetic.

Both argmaxes with certified structure are triple-binding, and their
architectures read
as follows. At the harmonic instance $v=(1,2,3,4,5)$, the deepest hole
$c^{*}=(\tfrac{122}{1175},\tfrac{688}{1175},\tfrac{284}{1175},
\tfrac{80}{1175})$ has four binding pair forms:
$\phi^{45}=-\rho$ at $z=(0,0,0,-1)$, $\phi^{34}=-\rho$ at
$z=(0,1,0,0)$, $\phi^{35}=+\rho$ at $z=(0,0,0,0)$, and
$\phi^{25}=+\rho$ at $z=(0,1,1,1)$. The first three are the triple
$\{3,4,5\}$ with signs $(-,-,+)$; Corollary~\ref{cor:denom}(i) gives
$2(3\cdot4+3\cdot5+4\cdot5)\,\rho=5k_{34}+3k_{45}-4k_{35}
=5\cdot4+3\cdot4-4\cdot0=32$, i.e.
\[
\rho\;=\;\frac{32}{2\cdot 47}\;=\;\frac{16}{47},
\qquad 47=3\cdot4+3\cdot5+4\cdot5.
\]
The harmonic argmax set is richer than a single point, and the identity
is the reason. Along the branch region containing the argmax, the three
dips of the triple $\{3,4,5\}$ obey the weighted identity with
\emph{constant} total,
$32\,r_{(3,5)}+27\,r_{(4,5)}+35\,r_{(3,4)}=32$ (the $c$-terms cancel
in \eqref{eq:weights}), so at every point of the region the smallest
dip is at most $32/94=16/47$; on the one-parameter \emph{balance
family}
\[
c(k)\;=\;\tfrac{1}{235}\,(2k{+}160,\;3k{-}13,\;4k{+}91,\;5k),
\qquad k=20,\dots,35,
\]
the three dips tie at exactly $16/47$, so the bound is attained along a
whole band. The family member $c(35)=(\tfrac{46}{47},
\tfrac{92}{235},\tfrac{231}{235},\tfrac{35}{47})$ is re-asserted
exactly by the cross-validation run of Appendix~\ref{sec:prov}
alongside the isolated argmax of Table~\ref{tab:refute}; both are
components of the argmax set, and the two-sided certificate of
\S\ref{sec:bb} covers them together.
At $v=(1,2,3,4,5,12)$ the deepest hole
$c^{*}=(\tfrac{4714}{5875},\tfrac18,\tfrac58,\tfrac6{25},
\tfrac{2034}{5875})$ has three binding pair forms: $\phi^{2,12}=-\rho$
at $z=(1,0,0,0,-1)$, $\phi^{5,12}=+\rho$ at $z=(1,0,1,0,1)$, and
$\phi^{25}=+\rho$ at $z=(1,1,2,2,4)$---the triple $\{2,5,12\}$ with
signs $(+{,}+{,}-)$ in the $(ab),(bc),(ac)$ roles; the identity gives
$2(10+24+60)\,\rho=5k_{2,12}-12k_{25}-2k_{5,12}=70-12+10=68$, i.e.
\[
\rho\;=\;\frac{68}{2\cdot 94}\;=\;\frac{17}{47},
\qquad 94=2\cdot5+2\cdot12+5\cdot12.
\]
(For the harmonic instance the same substitution reads
$5k_{34}+3k_{45}-4k_{35}=32$.)
The weighted identity \eqref{eq:weights} is exactly the relation
$35\,\phi^{34}+27\,\phi^{45}=32\,\phi^{35}$ (resp.\
$84\,\phi^{25}+34\,\phi^{5,12}=70\,\phi^{2,12}$) that the binding
values satisfy; summed over the three binding pairs it reads
$35r_{34}+27r_{45}+32r_{35}=32$ (resp.\ $84r_{25}+34r_{5,12}+70r_{2,12}
=68$) with $r$ the common binding value, which is how the two argmax
values were first read off the computed deepest holes.

Three remarks sharpen the record. First, the lemma converts the deepest
hole problem at such points into arithmetic: once the binding triple
and its translates are known, $\rho$ is a one-line integer computation,
and the congruence class of the argmax
($\chi_{uv}\equiv s_{uv}(u{+}v)\rho \bmod 1$ for the binding pairs)
turns the exact search into a finite parametrized scan---the method
that located the $m=12$ argmax above. Second, the repeated denominator
$47$ is not a law: $ab+ac+bc=47$ for the triple $\{3,4,5\}$ and
$ab+ac+bc=94=2\cdot47$ for $\{2,5,12\}$, and the reductions by
$\gcd 2$ and $\gcd 4$ happen to land on the same $47$; at
$v=(1,2,3,4,5,7)$ the recorded witness $4/11$ is pair-bisecting
(Corollary~\ref{cor:denom}(ii) with the pair $(4,7)$:
$2\cdot11\,\rho=8$), a different archetype with denominator governed by
the pair sum $11$, and no $47$ appears---the denominator follows the
binding architecture. What remains true, and worth recording, is that
both computed triple-binding argmaxes share the $47$: a coincidence
the identity does not force (the two weight sums are $47$ and $94$, and
the divisors $2$ and $4$ are incidental) but which the two data points
do not demystify either. Whether the small-triple weight sums
$ab+ac+bc$ cluster at the scale where the deepest holes of small
instances live is a question the archive legitimately raises and does
not answer. Third, the
archetypes have different reach: triple binding with weights
$c(a{+}b)$, $a(b{+}c)$, $b(a{+}c)$ summing to $2(ab{+}ac{+}bc)$, and
pair bisection with weight $2(u{+}v)$---and between them they govern
every exact value in the archive. The four certified ladder values at
the $n=5$ multiples are pair-bisecting at the center class on the pair
$(1,m)$---verified by the same binding-structure extraction, the pair
binding at two translates with opposite signs at each of
$m=10,20,25,30$---with reduced denominators $22$, $42$, $13$, $62$ all
dividing $2(1{+}m)$; the two triple-binding argmax values divide
$2(ab{+}ac{+}bc)$. The boundary rows at exactly the
threshold (for example $m=6,7,8$ at $n=5$) bind neither archetype
cleanly and are certified at their boundary values directly.

\subsection{The deepest-hole certification ladder}

Table~\ref{tab:ladder} records the corrected modular law with its
certification levels. The law as originally stated (``deepest hole at
the center iff $n\mid m$'') is a small-$n$ phenomenon, genuine and exact
at $n=3$, certified at $n=4$, broken at $n=5$ by the harmonic instance
itself---now with the exact value $\rho=16/47$ against
$\delta=1/3$---and extinct at $n=6$; the record below is the corrected
version, and the rows improved by the branch-and-bound method are
marked.

\begin{table}[htbp]
\centering
\small
\caption{The deepest-hole ladder for $(1,\dots,n{-}1,m)$: when is the
maximum of $D$ attained at the center class ($\rho=\delta$)?
``Exact'' = complete exact proof; ``certified'' = complete proof with an
independent outer loop; ``consistent'' = search-level only.}
\label{tab:ladder}
\begin{tabularx}{\textwidth}{@{} c c >{\raggedright\arraybackslash}X @{}}
\toprule
$n$ & tested $m$ & status of ``deepest hole at center iff $n\mid m$''
\\
\midrule
$3$ & $3..24$ & exact: holds iff $3\mid m$; exact witnesses against all
others \\
$4$ & $4..16$ & certified: holds iff $4\mid m$ ($m=4,8,12$
away-certified; $m=16$ certified by branch and bound); exact witnesses
against all non-multiples \\
$5$ & $5..30$ & law false: $m=5$ (harmonic; closed
$\rho=16/47\ne1/3$) and $m=15$ fail exactly; $m=10,20,25,30$ certified
by branch and bound ($\rho=7/22$, $13/42$, $4/13$, $19/62$ exactly, at
the center class); all $20$ non-multiples fail exactly
\\
$6$ & $6..12,18,24$ & extinct: every tested $m$ fails exactly,
including every multiple of $6$ through $4n$; at $m=12$ the failure is
the exact witness $\rho\ge17/47>\delta=9/26$ \\
\bottomrule
\end{tabularx}
\end{table}

The $n=5$ multiples carry the hold signature (the torsion center
attains $\delta$; the $K=32$ grid maximum equals $\delta$ exactly; a
flat ridge passes through the center; a $288$-start Nelder--Mead search
with exact re-evaluation finds nothing above $\delta$), but at $d=4$ a
global exact certificate was beyond the program's earlier arrangement
machinery. (Only $m=10$ is a rung-2 instance; $m=20,25,30$ are the
higher multiples of $5$.) The branch-and-bound method of \S\ref{sec:bb}
closed all four. At $m=25$ and $m=30$ the trees empty in under five
minutes each ($10{,}414$ and $10{,}646$ certified leaves; $\rho=4/13$
and $\rho=19/62$ exactly), and at $m=20$ the tree empties in sixteen
minutes ($37{,}018$ leaves; $\rho=13/42$). At $m=10$ the first attempt
($25$ minutes, $143{,}228$ box evaluations, $61{,}038$ exact leaf
certificates) did not terminate: the unresolved set had total volume
below $6\times10^{-7}$, the exact depths at its sampled centers were at
most $1135/3584<7/22$, and the largest exact $U$-bounds among the
sampled survivors were all exactly $571/1792=7/22+9/19712$---a fixed
quantization offset, consistent with pure minimax loss rather than
genuine excess, though a residual volume below $6\times10^{-7}$ does not
exclude a deeper point hiding inside it. The record left two points
open to objection---the search evidence behind the label was thin, and
the residual reading was unproved---and both were answered in a renewed
pass that closed the row. On the search side: $4{,}194{,}304$ uniform
random points, the full $\{k/16\}^4$ and $\{k/32\}^4$ dyadic grids, a
$500{,}000$-point torsion neighborhood, $2{,}048$ batched random uphill
climbs ($1.8$ million evaluations), and $64$ Nelder--Mead polishes, all
screened on a fixed $s$-grid over $[0,r]\cup[1-r,1)$---a \emph{sound}
screen, because the grid minimum of $G(c,\cdot)$ is an upper bound for
$D(c)$, so no point above the threshold can escape it---followed by
exact rational re-evaluation of the $114$ deepest candidates: every
line of the search converges to the torsion center
$(\tfrac12,0,\tfrac12,\tfrac12)$, at which the depth is exactly $7/22$,
and nothing anywhere exceeds $7/22$. On the certificate side, the
renewed cascade adds the exact Lipschitz box test of \S\ref{sec:bb}:
at the first attempt's residual width ($1/512$ per side, half-width
$1/1024$) the test closes the observed gap exactly, since
$1135/3584+1/1024=0.31766\ldots<7/22$, and the re-run tree empties
($190{,}295$ box evaluations, $95{,}147$ splits, \emph{zero} survivors,
$1{,}286$ seconds on one core; $846$ of $846$ sampled float-stage
prunes re-confirmed in exact rational arithmetic; the harmonic negative
control still refuses; the $n=4$ certificates re-certify with identical
box counts). Together with the torsion center attaining $7/22$, this
certifies
\[
\rho(1,2,3,4,10)\;=\;\tfrac{7}{22}\quad\text{exactly},
\]
and since $7/22=\delta=1/2-\lambda$ with $\lambda=2/11$, the deepest
hole at $m=10$ is the center class itself. The $n=5$ multiples residue
now has a clean shape, recorded without forcing it into a law: the odd
multiples $k\in\{1,3\}$ fail exactly, the even and higher multiples
$k\in\{2,4,5,6\}$ are certified.

\subsection{The branch-and-bound certificate and its negative control}
\label{sec:bb}

The certification method is a second, independent construction, built on
the runner formulation \eqref{eq:runnerform} rather than on the
arrangement enumeration of the program's records. For a dyadic $c$-box
$C$ define
\[
U(C)\;=\;\min_{\sigma\in[0,1)}\;
\max\Bigl(\wnc{\sigma},\;\max_{2\le j\le n}\;
\sup_{c\in C}\wnc{c_j-v_j\sigma}\Bigr),
\]
the box-supremum analogue of \eqref{eq:runnerform} itself (the
$\wnc\sigma$ term is the depth of the first runner, whose coordinate is
fixed at $0$; the inner maximum runs over $j=2,\dots,n$). It is
computed exactly: each inner supremum is piecewise linear in $\sigma$
with a rational kink set, so the minimum of their maximum is attained at
a rational kink or crossing. By the minimax inequality,
$\max_C D\le U(C)$, so a box with exact $U(C)\le r$ is \emph{certified}
to contain no point of depth above $r$; boxes that fail are split
dyadically. If the tree empties, $\rho\le r$ is exactly certified by a
finite, auditable leaf list. Two soundness refinements matter in
practice: the $\sigma$-search is restricted to $[0,r]\cup[1-r,1)$, which
is lossless because any certifying $\sigma$ has $\wnc\sigma\le r$; and
equality $U(C)=r$ prunes, which the tight boxes need. One further
sound test completes the cascade: $D$ is $1$-Lipschitz in the sup norm
(for fixed $\sigma$ every term $\wnc{c_j-v_j\sigma}$ moves by at most
$|c_j-c'_j|$ while the $\wnc\sigma$ term does not move, and minimizing
over $\sigma$ preserves the bound), so a dyadic box $C$ with center
$c_0$ and half-width $h$ satisfies $\max_C D\le D(c_0)+h$ with $D(c_0)$
exactly computable in rational arithmetic; and before the minimax bound
$U(C)$ is attempted, the box's center depth is bounded above by the
minimum of $G(c_0,\cdot)$ over a fixed float grid on
$[0,r]\cup[1-r,1)$ (the grid minimum is an upper bound for $D(c_0)$),
with a one-in-sixty-four sample of such prunes re-confirmed exactly.
The cascade is then: the cheap grid test; the grid-Lipschitz test; the
exact Lipschitz box test; the exact minimax test $U(C)\le r$; split.

The method carries its own negative control, and the control behaves
exactly as it must. Run on the harmonic instance $v=(1,2,3,4,5)$ at
$r=\delta=\tfrac13$---where the truth is $\rho=16/47>\tfrac13$---the
branch and bound \emph{refuses to certify}: its survivor set persists
around the true deep holes, and the exact depths at the survivors'
centers reach $607/1792\approx0.3388>\tfrac13$, at
$c=(\tfrac{11}{256},\tfrac{163}{256},\tfrac{15}{256},\tfrac{79}{256})$.
The machinery thus demonstrates both halves of its own soundness profile
on one instance: it certifies only what is true, and it localizes the
false. The same control distinguishes the two $n=5$ worlds sharply: for
the harmonic instance the surviving excess stays above the threshold by
a fixed margin, while for the certified $n=5$ multiples
($m=10,20,25,30$) the surviving excess is pure minimax loss and decays
with the box size.

The argmax-certified value of Table~\ref{tab:refute} ($m=12$) is
certified as follows. The exact arrangement analysis is run on the
argmax neighborhood (a box of half-width $1/125$ around the located
deepest hole): every vertex of the induced hyperplane arrangement
inside the box, every ridge candidate, and every box corner is
evaluated exactly, and the maximum over the whole box is the displayed
value, attained at the recorded point. This certifies that no point of
the neighborhood is deeper than the witness. The two-sided close-out at
this instance---pruned leaves outside, center-plus-slack Lipschitz net
and exact region analyses over the survivors inside---did not
terminate within the stated budgets, and the row stands as an
argmax-certified strict refutation with global equality open.

The structural reason the close-out resists is worth recording, because
it delimits the method itself. The $m=12$ argmax is non-dyadic
($5875=5^3\cdot47$), strictly interior to every dyadic box at every
depth, so the box-supremum $U(C)$ exceeds $r$ strictly for every box
containing the argmax and the plain pruning rule never fires there;
the instances the plain cascade certifies ($m=16$, $20$, $25$, $30$,
and $10$) have dyadic (torsion) argmaxes, which is why their trees
empty. This is an obstruction to the \emph{plain} cascade, not to the
method: a closure at a non-dyadic argmax must let the certificate
follow the argmax's rational structure, and the harmonic instance
shows that such certificates exist. Two were added. The
\emph{sliding-$\sigma$ gap certificate}: for a fixed integer lift
vector $m$, a point $c$ admits a $\sigma$ with every runner term $\le
r$ iff the interval
$[\max_j (c_j-m_j-r)/v_j,\ \min_j (c_j-m_j+r)/v_j]$ intersects
$[0,r]\cup[1-r,1]$; each endpoint is affine in $c$ and the gap is
concave, so its minimum over a box sits at a vertex---$2^d$ vertex
checks certify the whole box, with $\sigma$ allowed to \emph{slide}
with $c$ along the band, which the single-$\sigma$ minimax bound
cannot do. The \emph{triple-identity certificate} is the weight
identity of \S\ref{sec:weights} in box form: exhibit three pair-dips of
a runner triple with weights $w_{(a,b)}=v_c(v_a{+}v_b)$ whose weighted
values sum, over the box, to a constant $S\le(\sum w)\,r$ (the
$c$-coefficients cancel exactly, verified symbolically per candidate);
at every $c$ the minimum dip is then $\le S/\sum w\le r$, and the
non-pair terms at each dip are checked $\le r$ by exact affine interval
arithmetic. Splitting is lattice-aligned at the dip lattice $q=470$:
dyadic midpoints alone can never align with the balance family's
rational slice edges, and boxes straddling them unpruned are exactly
what strands the plain cascade. With these, the tree at $r=16/47$
empties: $782{,}255$ box evaluations ($3{,}721$ cheap-grid,
$45{,}598$ grid-Lipschitz, $61{,}191$ gap, $197{,}330$ triple-identity,
$83{,}288$ exact-minimax prunes; $391{,}127$ splits; maximum depth
$81$), \emph{zero} survivors, $5{,}935$~s on one core; $4{,}841$
sampled prunes re-confirmed exactly. Combined with either component of
the argmax set,
\[
\rho(1,2,3,4,5)\;=\;\tfrac{16}{47}\quad\text{exactly.}
\]
The soundness profile matches the rest of the method: adversarial boxes
around known deeper holes refuse the new certificates at their
thresholds, and the $n=4$ and $m=10$ certificates are untouched.

The harmonic sextuple $v=(1,2,3,4,5,6)$ was run at its witness level
$29/80$ and its tree does not empty within the same budget class; its
row in Table~\ref{tab:refute} therefore remains a strict witness, with
equality open, and the archive records the asymmetry honestly: at
$d=5$ the survivor regions are larger and more numerous, and the exact
region analysis at the harmonic sextuple's argmax does not terminate
within the stated budget.

\subsection{Summary of the archive}

The archive's content, stated as a single paragraph: the
covering-radius strengthening of the Lonely Runner statement is false
from $n=5$, measured exactly at the extremal instance
($\rho(1,2,3,4,5)=16/47$, certified from both sides) and falsified at
four further instances with exact witnesses ($29/80$ and $4/11$;
$17/47$ argmax-certified and $1793/5000$ at the two former boundary
rows, both upgraded, the superseded search records
$97/288\to607/1792\to4183/12288$ retained as provenance), boundary
rows at $n=5$ ($m=6,7,8$) at exactly the threshold, and a modular-law
extinction record at $n=6$; the denominators of the exact values are
forced by the pair-character identity (Lemma~\ref{lem:chi},
Corollary~\ref{cor:denom}), whose two archetypes---triple binding and
pair bisection---govern every exact value in the archive and explain,
and delimit, the recurring $47$; the equivalent
center-depth form survives untouched (it is $\lrc$); the deepest-hole
modular law is exact at $n=3$, certified at $n=4$ (both sides, all
tested $m$, including $m=16$ by the tree method), false at $n=5$ (the
$k=1,3$ multiples fail exactly, $k=1$ with the exact value
$\rho=16/47\ne\delta$; $k\in\{2,4,5,6\}$ are certified, each with
$\rho=\delta$ exactly), and extinct at $n=6$; and every number in the
paragraph is re-asserted by persisted, independent, exact computation.
The archive is finite, small, and open to inspection; that is the
point.
